\subsection{\texorpdfstring{同态 \(\mathbf{B} \otimes_{\mathbf{A}} \operatorname{Tor}^{\mathbf{A}}(\mathbf{E}, \mathbf{F}) \to \operatorname{Tor}^{\mathbf{B}}(\mathbf{E} \otimes_{\mathbf{A}} \mathbf{B}, \mathbf{B} \otimes_{\mathbf{A}} \mathbf{F})\) 与 \(\mathbf{B} \otimes_{\mathbf{A}} \operatorname{Ext}_{\mathbf{A}}(\mathbf{E}, \mathbf{F}) \to \operatorname{Ext}_{\mathbf{B}}(\mathbf{B} \otimes_{\mathbf{A}} \mathbf{E}, \mathbf{B} \otimes_{\mathbf{A}} \mathbf{F})\)}{Tor 与 Ext 的标量扩张同态}}

在本节中，我们假设 A 是交换的；给定一个环同态 \(\rho: A \to B\) 使得 \(\rho(A)\) 包含在 B 的中心内，并设 E 和 F 是两个 A-模。我们有 A-模复形的典范同构

\[
u: \mathrm{B} \otimes_ {\mathrm{A}} (\mathrm{L} (\mathrm{E}) \otimes_ {\mathrm{A}} \mathrm{L} (\mathrm{F})) \rightarrow (\mathrm{L} (\mathrm{E}) \otimes_ {\mathrm{A}} \mathrm{B}) \otimes_ {\mathrm{B}} (\mathrm{B} \otimes_ {\mathrm{A}} \mathrm{L} (\mathrm{F}));
\]

另一方面，由于 L(E) ⊗\_A B 和 B ⊗\_ L(F) 是自由 B-复形，则有分次 A-模的典范同态（X，第 102 页）

\[
\begin{array}{r l} \psi^ {\prime} (\mathrm{L(E)} \otimes_ {\mathrm{A}} \mathrm{B}, \mathrm{B} \otimes_ {\mathrm{A}} \mathrm{L(F)}): & \mathrm{H} ((\mathrm{L(E)} \otimes_ {\mathrm{A}} \mathrm{B}) \otimes_ {\mathrm{B}} (\mathrm{B} \otimes_ {\mathrm{A}} \mathrm{L(F)})) \\ & \to \operatorname{Tor} ^ {\mathrm{B}} (\mathrm{E} \otimes_ {\mathrm{A}} \mathrm{B}, \mathrm{B} \otimes_ {\mathrm{A}} \mathrm{F}) \end{array}
\]

最后，我们有一个同态（X，第 62 页）

\[
\gamma : \mathrm{B} \otimes_ {\mathbf {A}} \operatorname{Tor} ^ {\mathbf {A}} (\mathrm{E}, \mathrm{F}) \rightarrow \mathrm{H} (\mathrm{B} \otimes_ {\mathbf {A}} \mathrm{L} (\mathrm{E}) \otimes_ {\mathbf {A}} \mathrm{L} (\mathrm{F})) .
\]

分次 B-模的典范同态

\[
\mathrm{B} \otimes_ {\mathrm{A}} \operatorname{Tor} ^ {\mathrm{A}} (\mathrm{E}, \mathrm{F}) \rightarrow \operatorname{Tor} ^ {\mathrm{B}} (\mathrm{E} \otimes_ {\mathrm{A}} \mathrm{B}, \mathrm{B} \otimes_ {\mathrm{A}} \mathrm{F})\tag{13}
\]

定义为复合 \(\psi^{\prime}(\mathrm{L}(\mathrm{E})\otimes_{\mathrm{A}}\mathrm{B},\mathrm{B}\otimes_{\mathrm{A}}\mathrm{L}(\mathrm{F}))\circ \mathrm{H}(u)\circ \gamma .\)

命题 9. --- 若 B 在 A 上是平坦的，则同态 (13) 是双射的。

事实上， \(\psi^{\prime}(\mathbf{L}(\mathbf{E})\otimes_{\mathbf{A}}\mathbf{B},\mathbf{B}\otimes_{\mathbf{A}}\mathbf{L}(\mathbf{F}))\) 是双射的（X，第 102 页，命题 1），且 \(\gamma\) 是双射的（X，第 66 页，推论 2）。

假设 B 在 A 上是平坦的。在同态 (12) 中将 E 代入 Q，B 代入 N，并将 \(\mathbf{B} \otimes_{\mathbf{A}} \mathbf{F}\) 代入 M，我们得到一个同态

\[
\operatorname{Ext} _ {\mathbf {A}} \left(\mathrm{E}, \mathrm{B} \otimes_ {\mathbf {A}} \mathrm{F}\right)\rightarrow \operatorname{Ext} _ {\mathbf {B}} \left(\mathrm{B} \otimes_ {\mathbf {A}} \mathrm{E}, \mathrm{B} \otimes_ {\mathbf {A}} \mathrm{F}\right)
\]

由命题 8 可知它是双射的。在同态 (10) 中将 E 代入 M，F 代入 N，A 代入 B，B 代入 Q，并交换张量积的因子，我们得到一个 B-模的同态

\[
\mathrm{B} \otimes_ {\mathrm{A}} \operatorname{Ext} _ {\mathrm{A}} (\mathrm{E}, \mathrm{F}) \rightarrow \operatorname{Ext} _ {\mathrm{A}} (\mathrm{E}, \mathrm{B} \otimes_ {\mathrm{A}} \mathrm{F}),
\]

由此通过复合得到一个同态，称为典范同态

\[
\mathbf {B} \otimes_ {\mathbf {A}} \operatorname{Ext} _ {\mathbf {A}} (\mathbf {E}, \mathbf {F}) \rightarrow \operatorname{Ext} _ {\mathbf {B}} (\mathbf {B} \otimes_ {\mathbf {A}} \mathbf {E}, \mathbf {B} \otimes_ {\mathbf {A}} \mathbf {F}).\tag{14}
\]

命题 10. --- 同态 (14) 在以下情形中是双射的：

\begin{enumerate}
\def\labelenumi{\alph{enumi})}
\item
  B 是有限生成投射 A-模；
\item
  B 是平坦 A-模，A 是诺特的，且 E 是有限生成 A-模。这由命题 7（X，第 108 页）得出。
\end{enumerate}

